\subsection{Multiple-Choice Test}

The following multiple-choice test was administered to participants. Each question targets a specific programming concept, with one correct answer per question.

% Question 1
\begin{enumerate}

	\item \textbf{Modularity} \\
	      You are building a game with 5 enemies that all move in the same way. \\[0.5em]
	      What is the best way to structure the program?
	      \begin{enumerate}
		      \item Copy the same movement logic for each enemy.
		      \item Create one enemy behaviour and reuse it for all enemies.
		      \item Put all enemy logic into one very large block of code.
		      \item I don't know.
	      \end{enumerate}
	      \textit{Correct answer: B}

	      % Question 2
	\item \textbf{Event Handling} \\
	      In a game, the player collects a coin. The score should increase immediately. \\[0.5em]
	      What should trigger the score update?
	      \begin{enumerate}
		      \item A time interval that increases the score by the number of coins touched in the previous interval.
		      \item An event that occurs when the player touches the coin.
		      \item A random event in the program that eventually fires when moving close to a coin.
		      \item I don't know.
	      \end{enumerate}
	      \textit{Correct answer: B}

	      % Question 3
	\item \textbf{Data / State} \\
	      A game has a score variable that increases when coins are collected. \\[0.5em]
	      Which approach best organises the score data?
	      \begin{enumerate}
		      \item Each object has its own score variable to avoid mixing up coin scores.
		      \item The score is written directly into multiple parts of the program wherever the score is needed.
		      \item The score is stored in one place and updated when coins are collected.
		      \item I don't know.
	      \end{enumerate}
	      \textit{Correct answer: C}

	      % Question 4
	\item \textbf{State Machines} \\
	      A game character should be able to be in the states: running, jumping, and dead. The character switches state based on input and in-game events. \\[0.5em]
	      What is the best way to handle the character's states?
	      \begin{enumerate}
		      \item Create one state machine that gathers all states in one place and controls the transitions between them.
		      \item Use many if-statements spread throughout the code to check the character's situation wherever the behaviour needs to change.
		      \item Create a separate script for each state, each keeping track of when the other scripts should be activated.
		      \item I don't know.
	      \end{enumerate}
	      \textit{Correct answer: A}

	      % Question 5
	\item \textbf{Instantiation} \\
	      You are building a game with three types of coins: gold (10 points), silver (5 points), and bronze (1 point). You have created a ``Coin'' behaviour that handles collection and scoring. \\[0.5em]
	      What is the best way to create the three coin types?
	      \begin{enumerate}
		      \item Create three separate behaviours (GoldCoin, SilverCoin, BronzeCoin) with different point values.
		      \item Create one ``Coin'' behaviour with a point value variable, and then create three instances with different values.
		      \item Create one coin, copy it three times, and then use a script to change the coin values based on their position on the screen.
		      \item I don't know.
	      \end{enumerate}
	      \textit{Correct answer: B}

	      % Question 6
	\item \textbf{Instantiation} \\
	      You are making a game where enemies appear as the player shoots them down. \\[0.5em]
	      What happens when the program continuously creates new enemies and deletes them again?
	      \begin{enumerate}
		      \item The program slows down over time, because creating and deleting objects takes time.
		      \item The program becomes faster, because there are always only a few enemies at a time.
		      \item It has no effect on the program's speed.
		      \item I don't know.
	      \end{enumerate}
	      \textit{Correct answer: A}

	      % Question 7
	\item \textbf{Modularity} \\
	      You are building a game where the player can move, shoot, take damage, and collect items. All of this is currently written in a single script. \\[0.5em]
	      What is the best way to organise this?
	      \begin{enumerate}
		      \item Keep everything in one script so it is always easy to find.
		      \item Move all logic to the level or scene script to have one centralised place.
		      \item Split the logic into separate components (e.g.\ Movement, Combat, Inventory), each responsible for one area.
		      \item I don't know.
	      \end{enumerate}
	      \textit{Correct answer: C}

	      % Question 8
	\item \textbf{Data / State} \\
	      Your game character can be in one of three states: running, jumping, or attacking. You want to ensure that only one state is active at a time. \\[0.5em]
	      What is the best way to keep track of this?
	      \begin{enumerate}
		      \item Use a separate True/False variable for each state (e.g., isRunning, isJumping, isAttacking).
		      \item Determine the state each frame by checking the player's position and velocity.
		      \item Store the current state in one variable (e.g., playerState = ``running'') and update it when the state changes.
		      \item I don't know.
	      \end{enumerate}
	      \textit{Correct answer: C}

	      % Question 9
	\item \textbf{State Machines} \\
	      A door in a game can be closed, open, or locked. \\[0.5em]
	      What best describes what a state machine does in this case?
	      \begin{enumerate}
		      \item It keeps track of which state the door is in, and determines when it can switch state.
		      \item It draws the door on the screen in the correct colour depending on its appearance.
		      \item It stores the door's geographical position in the game world.
		      \item I don't know.
	      \end{enumerate}
	      \textit{Correct answer: A}

	      % Question 10
	\item \textbf{Event Handling} \\
	      A button in a game should start the game when the player clicks it. \\[0.5em]
	      What is the best way to make the button respond to a click?
	      \begin{enumerate}
		      \item Check every frame whether the mouse is over the button and whether it has been clicked, then perform the action.
		      \item Listen for a ``click'' event on the button, and perform the action when the event occurs.
		      \item Create a variable that counts how many times the button has been clicked.
		      \item I don't know.
	      \end{enumerate}
	      \textit{Correct answer: B}

\end{enumerate}